\documentclass[12pt,a4paper]{article}
\usepackage[utf8]{inputenc}
\usepackage[T1]{fontenc}
\usepackage[spanish,es-nodecimaldot]{babel}
\usepackage{geometry,amsmath,amssymb,booktabs,float}
\usepackage{tikz,pgfplots}
\usetikzlibrary{arrows.meta,positioning,babel}
\pgfplotsset{compat=1.18}
\usepackage[hidelinks]{hyperref}
\geometry{margin=2.5cm}
\title{Tarea 1\\\large Solución de ejercicios de redes y probabilidad}
\author{Jorge Alberto Suárez Saldaña}
\date{\today}
\begin{document}
\maketitle

\section{Red no dirigida (figura 1)}
La red consta de $N=6$ nodos y $M=10$ aristas:
\[
E=\{ab,ae,af,bc,bd,be,cd,ce,df,ef\}.
\]
\begin{figure}[H]\centering
\begin{tikzpicture}[every node/.style={circle,draw,fill=black,inner sep=3pt},scale=1]
\node[label=above:$a$] (a) at (0,1.3) {};
\node[label=above:$b$] (b) at (2.5,2.5) {};
\node[label=right:$c$] (c) at (4.5,1) {};
\node[label=above:$d$] (d) at (5.4,2.9) {};
\node[label=below:$e$] (e) at (2.1,-0.6) {};
\node[label=below:$f$] (f) at (5,-1.4) {};
\draw (a)--(b) (a)--(e) (a) to[bend right=32] (f) (b)--(c) (b)--(d) (b)--(e) (c)--(d) (c)--(e) (d)--(f) (e)--(f);
\end{tikzpicture}
\caption{Representación de la red no dirigida.}
\end{figure}
\subsection*{a) Matriz de adyacencia}
\[
A=\begin{pmatrix}
0&1&0&0&1&1\\
1&0&1&1&1&0\\
0&1&0&1&1&0\\
0&1&1&0&0&1\\
1&1&1&0&0&1\\
1&0&0&1&1&0
\end{pmatrix}.
\]
La matriz es simétrica, como corresponde a una red no dirigida.
\subsection*{b) Grado promedio}
Los grados de $(a,b,c,d,e,f)$ son $(3,4,3,3,4,3)$. Por el lema del apretón de manos,
\[
\boxed{\langle k\rangle=\frac{\sum_i k_i}{N}=\frac{2M}{N}=\frac{20}{6}=\frac{10}{3}\approx3.3333.}
\]
\subsection*{c) Densidad}
Para un grafo simple no dirigido sin lazos, hay $N(N-1)/2$ aristas posibles. Así,
\[
\boxed{D=\frac{2M}{N(N-1)}=\frac{20}{6\cdot5}=\frac23\approx0.6667.}
\]
\subsection*{d) Distribución de grados}
Hay cuatro vértices de grado $3$ y dos de grado $4$. La distribución empírica es $P(3)=4/6$ y $P(4)=2/6$.
\begin{figure}[H]\centering
\begin{tikzpicture}
\begin{axis}[ybar,bar width=18pt,width=.80\textwidth,height=6cm,xlabel={Grado $k$},ylabel={Número de nodos},xtick={0,1,2,3,4,5},ymin=0,ymax=5,ytick={0,1,2,3,4,5},nodes near coords,grid=major]
\addplot coordinates {(0,0)(1,0)(2,0)(3,4)(4,2)(5,0)};
\end{axis}
\end{tikzpicture}
\caption{Frecuencias de grados de la red no dirigida.}
\end{figure}

\section{Red dirigida (figura 2)}
La red tiene $N=6$ vértices. La figura contiene dos arcos paralelos $b\to c$. Contando ambos, se identifican $M=14$ arcos:
\[
\begin{aligned}
&a\to e,\ a\to f,\quad b\to a,\ b\to e,\ b\to c\ (\text{dos veces}),\\
&c\to b,\ c\to d,\quad d\to b,\ d\to f,\quad
 e\to b,\ e\to c,\quad f\to d,\ f\to e.
\end{aligned}
\]
\begin{figure}[H]\centering
\begin{tikzpicture}[>=Stealth,every node/.style={circle,draw,fill=black,inner sep=3pt},scale=1]
\node[label=above:$a$] (a) at (0,1.2) {};
\node[label=above:$b$] (b) at (2.6,2.5) {};
\node[label=right:$c$] (c) at (4.5,1.25) {};
\node[label=above:$d$] (d) at (5.5,3) {};
\node[label=below:$e$] (e) at (2.15,-.65) {};
\node[label=below:$f$] (f) at (5.15,-1.5) {};
\draw[->] (b)--(a);
\draw[->] (a)--(e);
\draw[->] (a) to[bend right=31] (f);
\draw[->] (b) to[bend right=22] (e);
\draw[->] (e)--(b);
\draw[->] (b)--(c);
\draw[->] (b) to[bend right=27] (c);
\draw[->] (c) to[bend right=24] (b);
\draw[->] (c)--(d);
\draw[->] (d)--(b);
\draw[->] (d)--(f);
\draw[->] (f) to[bend right=24] (d);
\draw[->] (e)--(c);
\draw[->] (f)--(e);
\end{tikzpicture}
\caption{Reconstrucción de la red dirigida con sus 14 arcos.}
\end{figure}
\subsection*{a) Matriz de adyacencia}
La entrada $(b,c)$ vale $2$ porque existen dos arcos desde $b$ hasta $c$:
\[
A=\begin{pmatrix}
0&0&0&0&1&1\\
1&0&2&0&1&0\\
0&1&0&1&0&0\\
0&1&0&0&0&1\\
0&1&1&0&0&0\\
0&0&0&1&1&0
\end{pmatrix}.
\]
\subsection*{b) Grado medio de entrada y de salida}
Las sumas de filas dan $k^{\mathrm{out}}=(2,4,2,2,2,2)$, y las de columnas dan $k^{\mathrm{in}}=(1,3,3,2,3,2)$. Cada arco suma una unidad a cada total, por lo que
\[
\boxed{\langle k^{\mathrm{in}}\rangle=\langle k^{\mathrm{out}}\rangle=\frac{M}{N}=\frac{14}{6}=\frac73\approx2.3333.}
\]
\subsection*{c) Densidad}
Una red dirigida simple sin lazos admite $N(N-1)=30$ pares ordenados distintos. Si se utiliza el cociente de \emph{arcos totales} respecto a esos pares (contando multiplicidad),
\[
\boxed{D_{\text{con multiplicidad}}=\frac{14}{30}=\frac7{15}\approx0.4667.}
\]
En cambio, la definición usual de densidad para un \emph{grafo dirigido simple} primero ignora los arcos paralelos. Como existen 13 pares ordenados conectados distintos,
\[
\boxed{D_{\text{grafo simple subyacente}}=\frac{13}{30}\approx0.4333.}
\]
Se incluyen ambos valores para precisar la convención utilizada.
\subsection*{d) Distribución de grados de entrada}
Las frecuencias son: un nodo con grado de entrada $1$, dos con grado $2$ y tres con grado $3$.
\begin{figure}[H]\centering
\begin{tikzpicture}\begin{axis}[ybar,bar width=18pt,width=.80\textwidth,height=5.8cm,xlabel={Grado de entrada $k^{\mathrm{in}}$},ylabel={Número de nodos},xtick={0,1,2,3,4},ymin=0,ymax=4,ytick={0,1,2,3,4},nodes near coords,grid=major]\addplot coordinates {(0,0)(1,1)(2,2)(3,3)(4,0)};\end{axis}\end{tikzpicture}
\caption{Distribución de grados de entrada.}
\end{figure}
\subsection*{e) Distribución de grados de salida}
Cinco vértices tienen grado de salida $2$ y uno tiene grado de salida $4$.
\begin{figure}[H]\centering
\begin{tikzpicture}\begin{axis}[ybar,bar width=18pt,width=.80\textwidth,height=5.8cm,xlabel={Grado de salida $k^{\mathrm{out}}$},ylabel={Número de nodos},xtick={0,1,2,3,4,5},ymin=0,ymax=6,ytick={0,1,2,3,4,5,6},nodes near coords,grid=major]\addplot coordinates {(0,0)(1,0)(2,5)(3,0)(4,1)(5,0)};\end{axis}\end{tikzpicture}
\caption{Distribución de grados de salida.}
\end{figure}

\section{Probabilidad de usuarios transmitiendo en una LAN}
Se comparte un enlace de $1\,\mathrm{Gbps}=1000\,\mathrm{Mbps}$, cada usuario activo requiere $50\,\mathrm{Mbps}$ y transmite el $15\%$ del tiempo. Supondremos independencia entre usuarios y probabilidades constantes de transmisión, como requiere el modelo binomial. El enlace permite $1000/50=20$ usuarios simultáneos sin exceder su capacidad nominal.
\subsection*{a) Probabilidad individual de transmisión}
\[\boxed{p=0.15=15\%.}\]
\subsection*{b) Probabilidad de exactamente $n$ usuarios activos}
Sea $X$ el número de usuarios transmitiendo de un total de 90. Entonces
\[
X\sim\operatorname{Binomial}(90,0.15),\qquad
\boxed{P(X=n)=\binom{90}{n}(0.15)^n(0.85)^{90-n}},
\]
para $n=0,1,\ldots,90$, donde
\[\binom{90}{n}=\frac{90!}{n!(90-n)!}.\]
Su esperanza es $E[X]=90(0.15)=13.5$ usuarios.
\subsection*{c) Probabilidad de al menos 21 usuarios activos}
Como 21 transmisiones requieren $21(50)=1050\,\mathrm{Mbps}>1000\,\mathrm{Mbps}$, se solicita la cola de la distribución binomial:
\[
\begin{aligned}
P(X\ge21)&=\sum_{n=21}^{90}\binom{90}{n}(0.15)^n(0.85)^{90-n}\\
&=1-\sum_{n=0}^{20}\binom{90}{n}(0.15)^n(0.85)^{90-n}\\
&\approx \boxed{0.0240055892\approx 0.0240=2.40\%.}
\end{aligned}
\]
Este resultado coincide con la respuesta aproximada indicada en el documento. Significa que bajo el modelo binomial, alrededor de $2.4\%$ de los instantes habrá al menos 21 usuarios transmitiendo, por lo que la demanda simultánea superaría $1\,\mathrm{Gbps}$.
\end{document}